
@@M0@@
(EGA @@M1@@ 14.3.2.1), d'où
@@M2@@
Puisque @@M3@@ est universellement caténaire, (EGA IV 5.6.5.1.) est applicable, et donne
@@M4@@
où @@M5@@ et @@M6@@ sont les deux termes à droite dans (1.5.4). Alors on conclut immédiatement de (1.5.4) que
@@M7@@

On voit donc que (1.3.2) et (1.3.3) restent valables pour un schéma noethérien universellement caténaire @@M8@@ de dimension @@M9@@, si l'on définit la filtration de @@M10@@ au moyen de (1.5.1) suivant la \emph{codimension} des supports. On notera que cette filtration coincide avec celle de (1.1.1), lorsque de plus @@M11@@ est \emph{biéqui-dimensionnel}.

\subsection*{1.6.}
On peut se demander s'il est toujours possible de définir sur un schéma noethérien @@M12@@ de dimension finie une fonction @@M13@@ possédant la propriété (i) de 1.4 et de plus la propriété suivante, plus forte que (ii) :
\begin{quote}
(ii bis) Si @@M14@@, et @@M15@@ est une spécialisation immédiate de @@M16@@, alors @@M17@@.
\end{quote}
